\section{Results}
\label{sec:results}

% Information -> estimator -> structure. All scored numbers live here,
% under the evaluation protocol and the methods of Section 4.

The results proceed from estimator comparisons to
nonlinear structure.
Sections~\ref{sec:dense_weak} and \ref{sec:dense_weak_ew} establish the
coefficient structure of the exogenous signal --- dense but weak ---
and the window weighting.
Sections~\ref{sec:tree_edge} and \ref{sec:three_block_results} measure
the nonlinear margin and close with the final three-block model.

\subsection{Lasso versus Ridge}
\label{sec:dense_weak}

The exogenous panel carries information beyond the HAR lag
structure of \citet{Corsi2009}; whether any of it is recovered is
decided by the estimator. This subsection holds the information set
fixed and varies only the penalty of the linear estimator. Two estimators are
compared: the lasso \citep{Tibshirani1996}, which applies an $\ell_1$ penalty
and sets a subset of coefficients to exactly zero, and ridge
\citep{HoerlKennard1970}, which applies an $\ell_2$ penalty and shrinks all
coefficients toward zero without setting any to zero. The result of the
comparison is a property of how the predictive coefficients are distributed
across the columns of the design. In this section that property is measured
directly: ridge attains lower out-of-sample QLIKE than the lasso on the
full design (Section~\ref{sec:dense_weak_headtohead}); the lasso zeroes
$304$ of $523$ exogenous columns on the full basis, of which $235$ are dead
or constant availability indicators
(Section~\ref{sec:dense_weak_diagnosis}); and five principal components
retain $72\%$ of the feature variance while losing essentially all of the
predictive signal (Section~\ref{sec:dense_weak_diagnosis}). The coefficient structure is
\emph{dense but weak}, in the sense of the sparse-versus-dense
prediction literature
\citep{GiannoneLenzaPrimiceri2021,KozakNagelSantosh2020,ShenXiu2025}:
many nonzero coefficients, none individually large. Section~\ref{sec:dense_weak_twoblock}
defines the estimator built for this structure: a two-block ridge with
penalty $\alpha = 1$ on the HAR-of-target block and $\alpha = 100$ on the
HAR-of-all-features block, with no selection step.

\paragraph{The head-to-head.}\label{sec:dense_weak_headtohead}
\paragraph{Design.}
The two penalized linear estimators are run on the full feature design ---
every exogenous column at once --- on one shared out-of-sample panel of
$n = 218{,}934$ thirty-minute bars, with identical rows in both cells;
nothing varies between the two arms except the penalty.

\paragraph{Comparison protocol.}
Both arms are compared against the shared incumbent of
Section~\ref{sec:data_prep}: per-bar OLS on HAR plus calendar features with
no exogenous inputs, at the benchmark QLIKE of \unifAzeroQLIKE{}. For each
arm, let $\ell^{\mathrm{arm}}_t$ and
$\ell^{\mathrm{inc}}_t$ be the per-bar QLIKE losses \citep{Patton2011} of
the arm and the incumbent at bar $t$, and let
$d_t = \ell^{\mathrm{arm}}_t - \ell^{\mathrm{inc}}_t$. The Diebold--Mariano
statistic \citep{DieboldMariano1995} is the sample mean $\bar d$ divided by
its Newey--West HAC standard error, with the small-sample correction of
\citet{HarveyLeybourneNewbold1997} applied; a negative $t$ favors the arm.
\decide{panel unification: this is the frozen causal-tuning battery on the
$218{,}934$-bar panel; the restructure plan flags that the paper's final
tables may need this re-run on the paper's one-panel convention}

% Source: drafts_2026-07-30/linear_vs_nonlinear.tex (linear columns of the
% battery), itself sourced from
% writeup/metrics_table_causal_tune_plus_spectral.csv (family == 'causal-tuned')
% cross-checked against writeup/causal_tune_battery_table.tex.
\paragraph{Result.}
The comparison is run on the full design --- the HAR ladder, the calendar
features, and the full exogenous block --- with both estimators under the
same causal tuning protocol, so that neither family carries an informed
penalty choice the other lacks
(Table~\ref{tab:lasso_vs_ridge}). Ridge
(\unifBOneRidgeTunedQLIKE{}, DM \unifBOneRidgeTunedDM{} against the
incumbent) beats the lasso (\unifBTwoLassoTunedQLIKE{},
\unifBTwoLassoTunedDM{}), and the elastic net --- the interpolating
family at the battery's mixed-penalty grid --- lands between the two
(\unifBThreeEnetTunedQLIKE{}, \unifBThreeEnetTunedDM{}): the loss is
monotone in how densely the penalty spreads the coefficients. Keeping
every column and shrinking produces a lower loss than selecting a subset,
when every family chooses its penalty by the same causal rule.

\begin{table}[H]
\centering
\caption{Ridge versus lasso on the full design, out-of-sample QLIKE with
the Diebold--Mariano $t$-statistic against the OLS--HAR incumbent in
parentheses. Panel: the frozen panel of
Section~\ref{sec:data_prep}, $n = 273{,}554$ identical rows in every
cell; rows scored under the first-iteration scalar protocol (incumbent
\unifAzeroQLIKEScalar{}), pending the one-protocol refresh (the lasso
arm is in harvest); negative $t$ favors the arm. Both arms re-select their penalty every $250$ solves from the
battery grids on a forward validation split.}
\label{tab:lasso_vs_ridge}
\begin{tabular}{@{}lccc@{}}
\toprule
 & Ridge & Elastic net & Lasso \\
\midrule
Causally tuned & \unifBOneRidgeTunedQLIKE{} (\unifBOneRidgeTunedDM{}) & \unifBThreeEnetTunedQLIKE{} (\unifBThreeEnetTunedDM{}) & \unifBTwoLassoTunedQLIKE{} (\unifBTwoLassoTunedDM{}) \\
\midrule
Incumbent (per-bar OLS--HAR) & \multicolumn{3}{c}{\unifAzeroQLIKEScalar{}} \\
\bottomrule
\end{tabular}
\end{table}

A penalty-level sweep locates where the tuned arms operate. Holding the
ridge family fixed and varying only the constant, QLIKE falls
monotonically in $\alpha$
across the tested grid: $0.23219$ at $\alpha = 0.1$, $0.23159$ at $0.3$,
$0.23083$ at $1$, $0.23003$ at $3$, and $0.22909$ at $10$, still
decreasing at the grid's edge; the tuned-ridge arm, which re-selects from
a grid reaching $\alpha = 1000$, lands below every fixed point. The
penalty \emph{level} moves the loss by more than the penalty \emph{family}
does at any fixed level, and the wide design rewards heavier shrinkage
than the $\alpha = 1$ default --- the direction the block penalties of
Section~\ref{sec:dense_weak_twoblock} take further. One caveat: all DM
statistics are computed against the incumbent, a third model, so the
ridge-minus-lasso differences themselves carry no significance statement
yet. \pending{pairwise ridge-vs-lasso DM (or the incumbent-cancelling
bound) on the full design, so the head-to-head margin is tested rather
than pointed at}

\paragraph{The diagnosis: dense, not sparse.}\label{sec:dense_weak_diagnosis}
The lasso's own solution path shows what $\ell_1$ selection does on this
design. On the full basis the lasso zeroes $304$ of $523$ exogenous columns.
Of those $304$, $235$ --- $77\%$ --- are dead or constant availability
indicators. Most of the lasso's selection therefore removes columns that
carry no signal. Removing such columns does not lower QLIKE relative to
ridge, because $\ell_2$ shrinkage makes the same columns small without
making them zero. Where the $\ell_1$ penalty zeroes live columns, QLIKE
rises: the margin by which ridge beats the lasso is carried by the live
columns the lasso zeroes, not by the dead ones.

A principal-component check measures whether the spread of coefficients
across columns is a low-dimensional structure in a rotated basis.
Truncating the design to five principal components retains $72\%$ of the
feature variance and loses essentially all of the predictive signal. The
signal therefore does not lie in the high-variance directions of the
feature covariance. The measured structure is: (i) no small subset of
columns concentrates the signal --- the only columns the lasso can zero
without raising QLIKE are dead or constant ones; and (ii) no small set of
principal directions concentrates it. The coefficient structure has many
small nonzero entries in the coordinate basis and is not low-rank in the
data's covariance geometry. This is the dense-but-weak structure of
\citet{GiannoneLenzaPrimiceri2021}, verified on the solution path
rather than inferred from the loss ordering.

The two-block ridge scores QLIKE \unifBlkTwoUserQLIKE{}, with
Diebold--Mariano \unifBlkTwoUserDM{} against the OLS--HAR incumbent. Under
the documented convention (penalties $1$/$3\times 10^{3}$, 250-day window)
the same structure scores \unifBlkTwoDocQLIKE{} at \unifBlkTwoDocDM{}.
These are the first scored runs of this exact specification.
\pending{DM of the two-block ridge against single-penalty ridge on
\texttt{all\_features} --- computed from the same per-bar losses at harvest}

\subsection{Time-Weighting the Window}
\label{sec:dense_weak_ew}

Every estimator in this paper weights the $24{,}000$ bars of its
training window
equally. A direct measurement of coefficient stability motivates asking
whether it should: on non-overlapping windows across the panel, the
backbone's coefficient vector is $94$--$97\%$ directionally invariant
over twenty-two years, while the exogenous block's loses roughly $40\%$
of its direction --- once, on a timescale of $62$--$250$ trading days
--- and then holds a plateau for two decades. There is no ongoing decay
to track, but the oldest bars in a two-year window do describe a
slightly different exogenous coefficient than the newest.

Exponentially down-weighting the window tests this. The gram and
cross-moment are accumulated with weight $w^{\,t-s}$ on bar $s$,
truncated to the same $24{,}000$ bars so that $w = 1$ reproduces the
flat window exactly, with the half-life $H = \log 2 / \log(1/w)$ the
only new quantity. Three variants were run. With $H$ \emph{fixed} at
half the window ($12{,}000$ bars), the estimator improves on the flat
window: QLIKE \unifBlkThreeEwFixedHTwelveKQLIKE{} against the
three-block ridge's \unifBlkThreeTunedQLIKE{}, a paired increment of
\unifIncrBlkThreeEwFixedHTwelveKDQ{} at DM
\unifIncrBlkThreeEwFixedHTwelveKDM{}. With $H$ fixed at $3{,}000$ bars
it is worse (\unifIncrBlkThreeEwFixedHThreeKDQ{} at
\unifIncrBlkThreeEwFixedHThreeKDM{}), and --- consistent with the penalty-selection
results above --- letting the causal tuner \emph{select} $H$ from a grid is worse
still (\unifIncrBlkThreeEwTunedDQ{} at \unifIncrBlkThreeEwTunedDM{}):
the selected half-life flickers across the entire grid from one retune
to the next, and the flicker costs several times what the right fixed
value earns. Mild fixed down-weighting of the oldest half-window helps;
selecting the forgetting rate on a $125$-bar validation tail
eliminates the gain. This is the same asymmetry the penalty comparison of
Section~\ref{sec:dense_weak_headtohead} measures, appearing on a
temporal dial.

We report a related negative result. The hypothesis that serial
correlation of the residuals explains why validation-selected penalties
so far exceed in-sample risk-optimal ones was tested and fails on
measurement: the out-of-sample one-step residuals are nearly white
(integrated autocorrelation time $1.88$, first-lag autocorrelation
$0.03$), because per-bar refitting leaves consecutive forecasts sharing
almost no estimation error. A correction built on that hypothesis is
inert by construction, and the discrepancy between in-sample and
validation-selected shrinkage remains attributed to the loss geometry
--- the QLIKE criterion on levels against the squared-error criterion
the estimator is risk-optimal for --- which this paper measures but
does not resolve.

One further negative result closes the linear discussion. A factor
summary of the panel cannot replace the exogenous block outright: a
ladder of two-block models --- the backbone plus $K$ rotated directions
of the exogenous set, $K$ from $5$ to the full live rank, under both
eigenvalue and frozen predictive orderings, each causally tuned ---
trails the matched two-block model on the raw columns by roughly
$2 \times 10^{-3}$ at every rung, \emph{including full rank}
(\unifIncrBlkTwoRotVarFullKTunedDQ{} at DM
\unifIncrBlkTwoRotVarFullKTunedDM{}). At full rank under a uniform block
penalty a rotation cannot change a ridge fit, so the gap is not a
property of the rotation: it traces to the construction having
re-standardized its inputs by frozen frame-window statistics where the
production path uses the rolling scaler, an identity verified directly
(rotating the production-scaled columns reproduces the raw block's
fitted values to $4 \times 10^{-15}$). Rolling standardization of the
inputs is therefore worth on the order of $2 \times 10^{-3}$ against
frozen standardization --- several times the size of most increments
this paper reports --- and it, not the choice of basis, ordering, or
width, is the largest single design decision in the factor machinery.

\subsection{Gradient-Boosted Trees}
\label{sec:tree_edge}

The linear results are now in place: the exogenous signal is dense and
weak (Section~\ref{sec:dense_weak}), and the two-block ridge is the
measured linear response (Section~\ref{sec:dense_weak_twoblock}). The
remaining question is what nonlinear structure adds. First,
gradient-boosted trees beat the best linear arm on every feature set,
by a small margin. Second, decompositions of the tree fits attribute
the bulk of that margin to additive-plus-pairwise structure, and the
pairwise part to products of columns the panel already contains, gated
on the session clock. Third, those products --- written down as
explicit columns and shrunk as a block of their own
(Section~\ref{sec:product_block}) --- close the paper's model ladder
in the final three-block ridge (Section~\ref{sec:three_block_results}).

\paragraph{The expert bank.} The twenty frozen configurations of
Section~\ref{sec:estimators} were walked per bar on the full panel and
scored under the evaluation protocol. Seventeen of the nineteen
completed experts beat the benchmark, at QLIKE $0.21796$--$0.22108$
against the benchmark's \unifAzeroQLIKE{} (DM $-11.3$ to $-7.6$); the
two seeded Sobol points --- included deliberately as unoptimized
controls --- lose to it ($0.23064$ and $0.23853$), so the Optuna
selection is worth $1$--$2 \times 10^{-2}$ over a random draw. One
expert's harvest is still completing; the bank aggregates below are
computed on the nineteen complete experts over their full common row
set ($263{,}782$ bars, 2003 onward), and the twentieth's inclusion is
expected to move them only marginally.

\paragraph{The bank beats its members.} The bank is aggregated two ways
on the same trailing-$125$-bar tail with a $25$-bar embargo, refreshed
every $250$ bars (the identical schedule to the linear tuners): a
\emph{selection} arm, which fills the next block from the
lowest-tail-loss expert, and a \emph{hedged} arm, which applies
exponential weights $\propto \exp(-\eta \sum \ell)$ with
$\eta = \sqrt{8 \ln K / L}$ --- the standard weighted-average
forecaster rate, no free constant. All tree rows are restricted to
2003 onward, since the menu was tuned on the pre-2003 dev span. The
selection arm scores \unifTreeTunedQLIKE{} (DM \unifTreeTunedDM{}
against the benchmark); the
hedged arm scores \unifTreeHedgeQLIKE{} (DM \unifTreeHedgeDM{}) ---
\textbf{better than every
individual expert, including the best of them} ($0.21796$), a paired
improvement of $-0.0046$ over the strongest member. The ordering
(selection $<$ best expert $<$ hedge, with hedge
\unifIncrTreeHedgeDQ{} better than
selection at DM \unifIncrTreeHedgeDM{}) is the paper's recurring tuning lesson in its
purest form: causal selection on a short tail is expensive, and
weighting is cheaper than selecting. The selection trajectory splits
nearly evenly across families ($296$ XGBoost to $262$ LightGBM
boundaries) with era-specific leans (LightGBM in 2009 and 2013, XGBoost
in 2018), and no expert dominates: the most-selected fills $22\%$ of
boundaries.

\begin{table}[H]
\centering
\caption{The tree expert bank under the evaluation protocol
(Section~\ref{sec:smearing_contract}). Individual experts are scored on
the full panel; the bank aggregates are computed on their common row set
($263{,}782$ bars, 2003 onward). DM $t$ is against the OLS--HAR
benchmark (\unifAzeroQLIKE{}); negative favours the arm.
\label{tab:tree_bank}}
\begin{tabular}{@{}lcc@{}}
\toprule
 & QLIKE & DM $t$ \\
\midrule
Best individual expert & $0.21796$ & $-11.3$ \\
Median completed expert & $0.21969$ & --- \\
Weakest Optuna-selected expert & $0.22108$ & --- \\
Sobol controls (2, unoptimized) & $0.23064$, $0.23853$ & --- \\
\midrule
\texttt{tree\_tuned} (causal best-expert selection) & \unifTreeTunedQLIKE{} & \unifTreeTunedDM{} \\
\texttt{tree\_hedge} (exponential weights) & \unifTreeHedgeQLIKE{} & \unifTreeHedgeDM{} \\
\bottomrule
\end{tabular}
\end{table}

\paragraph{Comparison discipline against the linear ladder.} The tree
rows above are $2003$+ by construction ($263{,}782$ bars); the linear
arms are scored
on $248{,}686$--$273{,}554$ bars under their own row rules. Pooled
levels across these row sets are not comparable (the campaign's row-set
discipline of Section~\ref{sec:smearing_contract}) --- the hedged
bank's \unifTreeHedgeQLIKE{} and the linear champion's
\unifBlkFourProdBbExogTrailSdQLIKE{} are separated by
$2 \times 10^{-5}$ \emph{on different row sets}, which is exactly why
the trees-versus-products verdict is reserved for the scorer's paired
increments on common rows, registered as the block-ladder harvest
completes (Appendix~\ref{app:protocol_note}). What is comparable as
stated: the hedged bank's \unifTreeHedgeDelta{} improvement over the
benchmark at
DM \unifTreeHedgeDM{}, and the bank-internal ordering above.

% ── AWAITING AUTHOR APPROVAL (2026-08-06) ────────────────────────────────
% The passage from here through "...takes the hint literally." is held
% pending the author's review: it carries three decompositions (fANOVA
% shares, the expressivity ladder, the product distillation) from earlier
% builds on non-comparable panels. Do not treat as final until the \decide
% below is resolved.
\decide{APPROVE OR REVISE this ``Where the nonlinearity lives'' passage
(through ``takes the hint literally''): all three decompositions are
carried over from earlier builds on non-comparable panels --- fANOVA
$55/9/36$ with HAR $96\%$; the expressivity ladder $0.12757 \to 0.12108
\to 0.12080 \to 0.12033$ vs XGBoost $0.12094$; the product distillation
$58\%/62\%$. Held from final until reviewed.}

\paragraph{Where the nonlinearity lives.} Three decompositions of the tree
fits were computed in an earlier phase of this project.

\emph{Interaction order.} A functional-ANOVA decomposition of a fitted
depth-eight boosted tree attributes approximately $55\%$ of its explainable
variance to additive terms, $9\%$ to pairwise interactions, and $36\%$ to a
higher-order remainder concentrated in no identifiable interaction. Within
the additive part, the HAR block accounts for $96\%$.

\emph{Saturation at pairwise order.} On the deployed residual-stack panel
of the earlier build ($n = 194{,}934$; levels not comparable to the
battery), an expressivity ladder holds the base model fixed and varies only
the correction stage. A purely additive correction scores $0.12757$.
Adding bilinear pairwise interactions scores $0.12108$. Nonlinear pairwise
shapes score $0.12080$. The deployed bagged additive-plus-pairwise stage (a
GA2M in the sense of \citealt{Lou2013}) scores $0.12033$. An
Optuna-searched unrestricted-depth XGBoost correction stage scores
$0.12094$, which is worse than fitting no correction at all ($0.12081$).
The ladder improves monotonically up to pairwise order; the unrestricted
stage does not improve on the pairwise stages.

\emph{Distillation into products.} On an earlier build of the same panel
family, a penalized linear base scores $0.12516$ and a residualized
boosted-tree correction improves it to $0.12414$. Adding to the base a
single hand-written interaction --- the HAR lags multiplied by a late-day
session dummy, six columns --- scores $0.12457$, which recovers $58\%$ of
the tree correction's improvement. The full set of HAR-by-six-regime
interactions scores $0.12453$, or $62\%$. The session dummy in the
recovered term marks the late-day bars at which the change of sign in
HAR's persistence coefficient is documented in
Section~\ref{sec:data_prep}.

The tree experiments therefore yield a specification statement:
\textbf{the panel's recoverable nonlinearity is a sparse set of pairwise
products of columns the design already contains}. The rest of this section
constructs that set of products as an explicit feature block.

\subsection{Three-Block Ridge Results}
\label{sec:three_block_results}

The campaign executor runs the three-block arm under two conventions. The
stated convention is Equation~\eqref{eq:three_block} at a $24{,}000$-bar
($500$-day) rolling fit window: penalties
$(\alpha_1, \alpha_2, \alpha_3) = (1, 100, 1000)$. The documented
convention pins $(1, 3\times10^{3}, 3\times10^{4})$ at a $12{,}000$-bar
($250$-day) window. Under the stated convention the three-block ridge
scores QLIKE \unifBlkThreeUserQLIKE{}, with Diebold--Mariano
\unifBlkThreeUserDM{} against the OLS--HAR incumbent; under the
documented convention it scores \unifBlkThreeDocQLIKE{} at
\unifBlkThreeDocDM{}. The product-block \emph{increment} --- the paired
per-bar comparison of the three-block against the two-block ridge on the
same rows --- is $\Delta\mathrm{QLIKE} = $ \unifIncrBlkThreeUserDQ{} at
DM \unifIncrBlkThreeUserDM{} under the stated convention, and
\unifIncrBlkThreeDocDQ{} at \unifIncrBlkThreeDocDM{} under the documented
one; under causal per-block tuning, \unifIncrBlkThreeTunedDQ{} at
\unifIncrBlkThreeTunedDM{}.

\begin{table}[H]
\centering
\caption{The block ladder under the evaluation protocol
(Section~\ref{sec:smearing_contract}), stated and documented conventions.
DM $t$ is against the OLS--HAR benchmark (\unifAzeroQLIKE{}); negative
favours the arm. The product-block increment is the paired comparison of
the three-block against the two-block ridge on the same rows.
\label{tab:block_ladder}}
\begin{tabular}{@{}lccc@{}}
\toprule
 & QLIKE & DM $t$ & $n$ \\
\midrule
OLS on HAR $+$ calendar (benchmark) & \unifAzeroQLIKE{} & --- & $273{,}554$ \\
Two-block ridge, stated $(1, 100)$ & \unifBlkTwoUserQLIKE{} & \unifBlkTwoUserDM{} & $272{,}971$ \\
Two-block ridge, documented $(1, 3{\times}10^{3})$ & \unifBlkTwoDocQLIKE{} & \unifBlkTwoDocDM{} & $272{,}971$ \\
Three-block ridge, stated $(1, 100, 1000)$ & \unifBlkThreeUserQLIKE{} & \unifBlkThreeUserDM{} & $248{,}280$ \\
Three-block ridge, documented $(1, 3{\times}10^{3}, 3{\times}10^{4})$ & \unifBlkThreeDocQLIKE{} & \unifBlkThreeDocDM{} & $272{,}971$ \\
\midrule
Product-block increment (stated) & \unifIncrBlkThreeUserDQ{} & \unifIncrBlkThreeUserDM{} & $248{,}280$ \\
Product-block increment (documented) & \unifIncrBlkThreeDocDQ{} & \unifIncrBlkThreeDocDM{} & $272{,}971$ \\
\bottomrule
\end{tabular}
\end{table}
\decide{whether the three-block table, once run, reports the tree arms
alongside it --- the natural closing comparison is the three-block ridge
against the causally tuned LightGBM on \texttt{all\_features} ($0.12604$),
i.e.\ whether naming the tree's structure recovers the tree's level within
the linear class}

The product block does not carry all of the tree's fit: the
functional-ANOVA remainder and the expressivity ladder's residual gap sit
above pairwise order, and the block contains pairwise products only. What
the block does carry runs in a linear estimator --- one solver, three
penalties, refit every bar.
